\movedsub{Does this reach past entailment scorers?}
\label{sec:beyond-nli}

Two of the three evaluators we test are built on the same entailment model, which invites the reading that we have diagnosed NLI rather than faithfulness measurement. Other paradigms exist. ARES \citep{saad-falcon-etal-2024} fine-tunes lightweight judges on synthetic data and calibrates them against a few hundred human annotations with prediction-powered inference. RAGVUE \citep{murugaraj-etal-2026} decomposes a RAG system into retrieval quality, answer relevance, claim-level faithfulness and judge calibration, and attaches an explanation to each. We ran neither, so what follows is an argument about what our measurements imply for them, not a test of them.

Three of our results do not depend on entailment. AlignScore is not an NLI model --- it is an alignment function trained across factual-consistency tasks --- and in the verbosity control it loses more than either entailment aggregate, falling from 74.5\% to 18.5\% on HotpotQA with its mean response to a falsified value dropping 0.494 to 0.127 (Section~\ref{sec:verbosity}). The structural bias therefore attaches to scalar aggregation over a multi-assertion answer, not to the entailment model underneath. Any scorer that returns one number for a five-assertion answer inherits it, including a fine-tuned judge.

Claim-level decomposition is the obvious response and is what RAGVUE prescribes. We measure it. Claim-min is the most length-robust of the three, giving up 36 detection points between the extreme generators against 54 and 46 for the others, and it holds 43.6\% detection where the others fall below 19\%. It is also the aggregate whose absolute scores collapse at that verbosity, to means of 0.234 and 0.076, so decomposition buys sensitivity and costs thresholdability. That trade is a property of taking a minimum, not of the model taking it, and it will appear in any claim-decomposed metric.

Calibration is a different axis from the one we measure, which is why ARES is complementary rather than an answer. Prediction-powered inference corrects a judge's bias against a human-labelled sample: it makes the reported number closer to what annotators would say on that distribution. Our two failure modes survive that correction. A judge can be unbiased on average and still fail to separate two systems at the sample sizes people use, which is resolving power (Appendix~\ref{app:power}), and it can be unbiased on average and still move by 0.054 when a grounded value is replaced with a false one, which is sensitivity to a known error. Neither is a calibration error, and neither would be visible in agreement with annotators on unperturbed answers.

What would test this properly is running these frameworks through the same probe. We expect the verbosity effect to reproduce for any scalar aggregate and the calibration machinery to be orthogonal to it, and we have not shown either.
